\begin{lemma}\label{lem:octic-doubled-cuspidal-RR}
Let $X/F$ be a modular curve with a rational cusp $c_0$,
and suppose that $\operatorname{Jac}(X)(F)$ has rank zero.
Let $v\mid2$ be unramified, with good reduction for $X$, and
suppose that all cusps split over an extension unramified at $v$.
If
\[
 l(2C)=1
\]
for every effective cuspidal divisor $C$ of degree $d$ on
the residue curve, then no $F$-rational noncuspidal effective
divisor of degree $d$ has entirely cuspidal reduction at $v$.
\end{lemma}
\begin{proof}
Suppose $D$ has such reduction $\overline C$, and lift it
\emph{\'etale}-locally to a cuspidal divisor $C$ over an
extension unramified at $v$. Both $[D-dc_0]$ and $[C-dc_0]$
are torsion, the former by rank zero and the latter by
Manin--Drinfeld. Hence $[D-C]$ is torsion in the reduction
kernel. Over an unramified dyadic field this kernel is killed
by $2$: multiplication by $2$ sends the first formal-group
filtration into the second, on which the formal logarithm
is injective. Thus $2D\sim2C$.
The distinct effective divisors $2D$ and $2C$ give
$l(2C)\geq2$. Specialization and base change then give
$l(2\overline C)\geq2$ over the residue field, a contradiction.
No rank assertion over the cusp-splitting extension is used.
\end{proof}

\begin{proposition}\label{prop:octic-no321}
There is no elliptic curve over a number field of degree $8$
with a subgroup isomorphic to $C_3\times C_{21}$.
\end{proposition}
\begin{proof}
The Weil pairing gives $F=\mathbb Q(\zeta_3)\subseteq K$,
so a proposed modular point gives an $F$-rational
noncuspidal effective divisor $D$ of degree $4$ on
$X=X_1(3,21)$. This curve has genus $25$, has good reduction
at the prime above $2$, and has a rational cusp over $F$.
Its cusps split over $\mathbb Q(\zeta_{21})$, unramified at $2$.
By \cite[Theorem~4.1]{DS}, its Jacobian has rank zero over $F$.

The cusp signature over $\mathbb F_4$ is $1^{24},3^8$.
The exact noncuspidal closed-place counts in degrees one
through four are
\[
 (0,0,0,144).
\]
They agree with the independent Frobenius computation on
$\operatorname{Jac}(X_1(63)/\langle22\rangle)$ after quadratic
constant-field extension. Thus $\overline D$ is either
entirely cuspidal or one noncuspidal quartic place.

For the first case there are exactly
\[
 \binom{27}{4}+24\cdot8=17742
\]
effective cuspidal degree-$4$ divisors.
The fresh computation \texttt{double\_cusp\_3\_21\_D4.m}
checks $l(2C)=1$ for all of them. It computes the common
$25$-dimensional space of holomorphic differentials and
the $104$ local vanishing-constraint matrices for
$2m c$, with $m\deg c\leq4$.
For each $C$, their combined matrix has rank $8$;
Riemann--Roch gives $l(2C)=9-8=1$.
The matrices are exported for an independent exhaustive
$\mathbb F_4$ replay, and ordinary Riemann--Roch independently
checks the initial cases and every exceptional case.
A separate native function-field execution,
\texttt{double\_cusp\_3\_21\_D4\_dt.m}, reproduces all the
matrices by direct residue lifting: for a local uniformizer $u$
and a fixed $\mathbb F_4$-linear residue-field lift, repeatedly
subtract the lift of the residue of $f$ and divide by $u$,
starting with $f=\omega/\omega_0$.
This gives the same vanishing spaces without Hensel expansion.
It checks every one of the $104$ individual local dimensions
with ordinary Riemann--Roch and again obtains $l(2C)=1$
for all $17742$ divisors. Both independently exported matrix
sets pass the exhaustive $\mathbb F_4$ replay.
Lemma~\ref{lem:octic-doubled-cuspidal-RR} excludes this branch.

For the second case put $Q_3=7Q$, $Q_7=3Q$, so
$Q=Q_3+5Q_7$. The $F$-defined automorphisms
\[
 (P,Q)\longmapsto
 (aP+bQ_3,\ cP+dQ_3+5uQ_7),
 \qquad
 \begin{pmatrix}a&b\\c&d\end{pmatrix}\in\operatorname{SL}_2(\mathbb F_3),
 \quad u\in(\mathbb Z/7\mathbb Z)^\times,
\]
give a group of order $72$, after identifying simultaneous
negation. They partition the $144$ quartic places into
two orbits. The complete key partition is checked on every
place, and one full orbit is independently compared with
the original moduli inverse, fixing the $\mathbb F_4$ base
embedding in every residue-field map.

The prime $13$ splits in $F$. The rational Hecke operator
\[
 A_{13}=T_{13}-1-13\langle13\rangle
\]
annihilates $J(F)$, and the rational base cusp satisfies
$T_{13}c_0=c_0+13\langle13\rangle c_0$.
Thus $A_{13}D$ is principal. It commutes with the displayed group:
$13\equiv1\pmod3$ and the generator action on the
$7$-part is central. It therefore suffices to test the two
representatives. The fresh exact computation
\texttt{sieve\_3\_21\_D4\_p2.m} finds $A_{13}(x)$
nonprincipal for both. At the first representative,
the fast $T_{13}$ divisor agrees exactly with the
independent generic-isogeny construction and has degree $56$.
This contradicts the necessary principality of
$A_{13}(\overline D)$, and proves the proposition.
\end{proof}

